\chapter{状态矩阵}

\footnotesize
\begin{ReportWideTable}{>{\raggedright\arraybackslash}p{3.8cm}>{\raggedright\arraybackslash}p{3.5cm}>{\raggedright\arraybackslash}p{6.1cm}}
\captionof{table}{当前研究状态矩阵}\label{tab:status-matrix}
\toprule
\textbf{主题} & \textbf{研究状态} & \textbf{说明}\\
\midrule
Nous Wave 定位 & 目标设计 & 可组合的主体认知与人格平台；记忆型能力组合可独立使用。\\
Nous Wave / Heptalogos & 目标设计 & 共享 SubjectId；认知权威与行为、规范权威分别维护。\\
记忆（Memory） & 必需能力 & 基础能力组合中的认知领域所有者。\\
自我 / 关系 / 动机 & 高级目标 & 各领域合同单独确定，可与基础记忆能力组合启用。\\
人格风格（Persona） & 目标设计 & 自我领域中的一种表示。\\
认知图式（CognitiveSchema） & 必需能力；归属待定 & 表达跨经历形成的结构认知；具体领域归属待定。\\
联想拓扑（AssociativeTopology） & 目标设计 & 可重建的检索服务投影；认知权威由领域数据维护。\\
类型化有界 PPR/扩散 & 候选基线 & 作为可解释的拓扑对照方案。\\
Wave & 当前已实现；候选方案 & 当前代码默认使用 Wave，目标研究要求与其他基线比较。\\
EPA / Residual & 候选方案及局部证据 & 数学性质、局部试验和端到端效用分别评价。\\
稠密检索后端 & 候选方案 & Qdrant 与 pgvector 按本机工作负载评测。\\
自我构成 / 动机细节 & 待定问题 & 后续合同将明确领域归属、范围和写入规则。\\
\bottomrule
\end{ReportWideTable}
\normalsize
